% =====================================================================
%  Thème 3 — Colinéarité et parallélisme — Niveau DIFFICILE — Exercices
% =====================================================================
\documentclass[11pt,a4paper]{article}
\usepackage[utf8]{inputenc}
\usepackage[T1]{fontenc}
\usepackage{lmodern}
\usepackage[french]{babel}
\usepackage{amsmath,amssymb,mathtools}
\usepackage{xcolor}
\usepackage{enumitem}
\usepackage{fancyhdr}
\usepackage[a4paper,margin=2.1cm,headheight=26pt,headsep=10pt]{geometry}

\definecolor{primary}{RGB}{222,70,80}
\definecolor{primarydark}{RGB}{150,30,42}

\pagestyle{fancy}\fancyhf{}
\fancyhead[L]{\small\textcolor{primarydark}{\textbf{Fiche 2} — Calcul vectoriel}}
\fancyhead[R]{\small\textcolor{primarydark}{\textsc{Thème 3 — Difficile — Exercices}}}
\fancyfoot[C]{\small\thepage}
\renewcommand{\headrulewidth}{0.8pt}

\newcommand{\vc}[1]{\overrightarrow{#1}}
\providecommand{\norm}[1]{\left\lVert #1\right\rVert}
\newcommand{\exo}[1]{\medskip\noindent{\color{primary}\bfseries\large Exercice #1.}\par\nopagebreak\smallskip}

\begin{document}
\begin{center}
{\LARGE\bfseries\color{primarydark}Colinéarité et parallélisme — Niveau Difficile}\\[2pt]
{\color{primary}\rule{0.5\linewidth}{1.1pt}}
\end{center}
\vspace{1mm}
{\small\itshape Objectif : colinéarité avec des radicaux, résolution d'une condition paramétrée,
alignement dans une situation concrète et synthèse « vrai / faux » de type concours.}
\vspace{2mm}

\exo{1} \textbf{Colinéarité avec des radicaux.}
On pose $\vec u\,\bigl(\sqrt3\,;\,1\bigr)$.
\begin{enumerate}[label=\textbf{\alph*)},leftmargin=2em,itemsep=3pt]
  \item Montrer, par le déterminant, que $\vec v\,\bigl(3\,;\,\sqrt3\bigr)$ est colinéaire à $\vec u$,
        puis déterminer le réel $k$ tel que $\vec v=k\,\vec u$.
  \item Le vecteur $\vec w\,\bigl(2\,;\,2\sqrt3\bigr)$ est-il colinéaire à $\vec u$ ? Justifier.
\end{enumerate}

\exo{2} \textbf{Condition de colinéarité paramétrée.}
On considère $\vec u\,(m;2)$ et $\vec v\,(8;m)$, où $m$ est un réel.
\begin{enumerate}[label=\textbf{\alph*)},leftmargin=2em,itemsep=3pt]
  \item Exprimer le déterminant de $\vec u$ et $\vec v$ en fonction de $m$.
  \item Déterminer toutes les valeurs de $m$ pour lesquelles $\vec u$ et $\vec v$ sont colinéaires.
  \item Pour chaque valeur trouvée, préciser le réel $k$ tel que $\vec v=k\,\vec u$.
\end{enumerate}

\exo{3} \textbf{Mise en situation — le réseau d'antennes.}
Sur une carte (repère orthonormé, en kilomètres), trois antennes sont placées en
$A\,(1;1)$, $B\,(4;3)$ et $C\,(10;7)$.
\begin{enumerate}[label=\textbf{\alph*)},leftmargin=2em,itemsep=3pt]
  \item Montrer que les trois antennes $A$, $B$, $C$ sont alignées.
  \item On veut installer une nouvelle antenne $D$ \emph{sur la même droite}, d'abscisse $x_D=7$.
        Déterminer son ordonnée $y_D$.
  \item Vérifier que $\vc{BD}$ est bien colinéaire à $\vc{AB}$.
\end{enumerate}

\exo{4} \textbf{Synthèse « vrai / faux » (type concours).}
Dans un repère orthonormé, on donne $A\,(0;3)$, $B\,(1;5)$, $C\,(-2;-1)$ et la droite
$d:\ 4x-2y+5=0$. Pour chaque affirmation, répondre par \textbf{Vrai} ou \textbf{Faux} en justifiant.
\begin{enumerate}[label=\textbf{\Alph*.},leftmargin=2.2em,itemsep=3pt]
  \item Les points $A$, $B$, $C$ sont alignés.
  \item Le vecteur $\vec u\,(1;2)$ est un vecteur directeur de la droite $(AB)$.
  \item Les droites $(AB)$ et $d$ sont parallèles.
  \item Les droites $(AB)$ et $d$ sont confondues.
\end{enumerate}

\end{document}
